% Options for packages loaded elsewhere
\PassOptionsToPackage{unicode}{hyperref}
\PassOptionsToPackage{hyphens}{url}
\PassOptionsToPackage{dvipsnames,svgnames,x11names}{xcolor}
\documentclass[
  11pt,
  a4paper,
]{article}
\usepackage{xcolor}
\usepackage[margin=1in]{geometry}
\usepackage{amsmath,amssymb}
\setcounter{secnumdepth}{-\maxdimen} % remove section numbering
\usepackage{iftex}
\ifPDFTeX
  \usepackage[T1]{fontenc}
  \usepackage[utf8]{inputenc}
  \usepackage{textcomp} % provide euro and other symbols
\else % if luatex or xetex
  \usepackage{unicode-math} % this also loads fontspec
  \defaultfontfeatures{Scale=MatchLowercase}
  \defaultfontfeatures[\rmfamily]{Ligatures=TeX,Scale=1}
\fi
\usepackage{lmodern}
\ifPDFTeX\else
  % xetex/luatex font selection
  \ifLuaTeX
    \usepackage{luaotfload}
    \directlua{luaotfload.add_fallback("mainfontfallback",{
      "NotoSerifBengali:mode=harf"
    })}
  \fi
  \setmainfont[,RawFeature={fallback=mainfontfallback}]{Latin Modern Roman}
  \ifLuaTeX
    \usepackage{luaotfload}
    \directlua{luaotfload.add_fallback("sansfontfallback",{
      "NotoSansBengali:mode=harf"
    })}
  \fi
  \setsansfont[,RawFeature={fallback=sansfontfallback}]{Latin Modern Sans}
  \ifLuaTeX
    \usepackage{luaotfload}
    \directlua{luaotfload.add_fallback("monofontfallback",{
      "NotoSansBengali:mode=harf"
    })}
  \fi
  \setmonofont[,RawFeature={fallback=monofontfallback}]{DejaVu Sans Mono}
\fi
% Use upquote if available, for straight quotes in verbatim environments
\IfFileExists{upquote.sty}{\usepackage{upquote}}{}
\IfFileExists{microtype.sty}{% use microtype if available
  \usepackage[]{microtype}
  \UseMicrotypeSet[protrusion]{basicmath} % disable protrusion for tt fonts
}{}
\makeatletter
\@ifundefined{KOMAClassName}{% if non-KOMA class
  \IfFileExists{parskip.sty}{%
    \usepackage{parskip}
  }{% else
    \setlength{\parindent}{0pt}
    \setlength{\parskip}{6pt plus 2pt minus 1pt}}
}{% if KOMA class
  \KOMAoptions{parskip=half}}
\makeatother
\usepackage{color}
\usepackage{fancyvrb}
\newcommand{\VerbBar}{|}
\newcommand{\VERB}{\Verb[commandchars=\\\{\}]}
\DefineVerbatimEnvironment{Highlighting}{Verbatim}{commandchars=\\\{\}}
% Add ',fontsize=\small' for more characters per line
\usepackage{framed}
\definecolor{shadecolor}{RGB}{248,248,248}
\newenvironment{Shaded}{\begin{snugshade}}{\end{snugshade}}
\newcommand{\AlertTok}[1]{\textcolor[rgb]{0.94,0.16,0.16}{#1}}
\newcommand{\AnnotationTok}[1]{\textcolor[rgb]{0.56,0.35,0.01}{\textbf{\textit{#1}}}}
\newcommand{\AttributeTok}[1]{\textcolor[rgb]{0.13,0.29,0.53}{#1}}
\newcommand{\BaseNTok}[1]{\textcolor[rgb]{0.00,0.00,0.81}{#1}}
\newcommand{\BuiltInTok}[1]{#1}
\newcommand{\CharTok}[1]{\textcolor[rgb]{0.31,0.60,0.02}{#1}}
\newcommand{\CommentTok}[1]{\textcolor[rgb]{0.56,0.35,0.01}{\textit{#1}}}
\newcommand{\CommentVarTok}[1]{\textcolor[rgb]{0.56,0.35,0.01}{\textbf{\textit{#1}}}}
\newcommand{\ConstantTok}[1]{\textcolor[rgb]{0.56,0.35,0.01}{#1}}
\newcommand{\ControlFlowTok}[1]{\textcolor[rgb]{0.13,0.29,0.53}{\textbf{#1}}}
\newcommand{\DataTypeTok}[1]{\textcolor[rgb]{0.13,0.29,0.53}{#1}}
\newcommand{\DecValTok}[1]{\textcolor[rgb]{0.00,0.00,0.81}{#1}}
\newcommand{\DocumentationTok}[1]{\textcolor[rgb]{0.56,0.35,0.01}{\textbf{\textit{#1}}}}
\newcommand{\ErrorTok}[1]{\textcolor[rgb]{0.64,0.00,0.00}{\textbf{#1}}}
\newcommand{\ExtensionTok}[1]{#1}
\newcommand{\FloatTok}[1]{\textcolor[rgb]{0.00,0.00,0.81}{#1}}
\newcommand{\FunctionTok}[1]{\textcolor[rgb]{0.13,0.29,0.53}{\textbf{#1}}}
\newcommand{\ImportTok}[1]{#1}
\newcommand{\InformationTok}[1]{\textcolor[rgb]{0.56,0.35,0.01}{\textbf{\textit{#1}}}}
\newcommand{\KeywordTok}[1]{\textcolor[rgb]{0.13,0.29,0.53}{\textbf{#1}}}
\newcommand{\NormalTok}[1]{#1}
\newcommand{\OperatorTok}[1]{\textcolor[rgb]{0.81,0.36,0.00}{\textbf{#1}}}
\newcommand{\OtherTok}[1]{\textcolor[rgb]{0.56,0.35,0.01}{#1}}
\newcommand{\PreprocessorTok}[1]{\textcolor[rgb]{0.56,0.35,0.01}{\textit{#1}}}
\newcommand{\RegionMarkerTok}[1]{#1}
\newcommand{\SpecialCharTok}[1]{\textcolor[rgb]{0.81,0.36,0.00}{\textbf{#1}}}
\newcommand{\SpecialStringTok}[1]{\textcolor[rgb]{0.31,0.60,0.02}{#1}}
\newcommand{\StringTok}[1]{\textcolor[rgb]{0.31,0.60,0.02}{#1}}
\newcommand{\VariableTok}[1]{\textcolor[rgb]{0.00,0.00,0.00}{#1}}
\newcommand{\VerbatimStringTok}[1]{\textcolor[rgb]{0.31,0.60,0.02}{#1}}
\newcommand{\WarningTok}[1]{\textcolor[rgb]{0.56,0.35,0.01}{\textbf{\textit{#1}}}}
\usepackage{longtable,booktabs,array}
\usepackage{calc} % for calculating minipage widths
% Correct order of tables after \paragraph or \subparagraph
\usepackage{etoolbox}
\makeatletter
\patchcmd\longtable{\par}{\if@noskipsec\mbox{}\fi\par}{}{}
\makeatother
% Allow footnotes in longtable head/foot
\IfFileExists{footnotehyper.sty}{\usepackage{footnotehyper}}{\usepackage{footnote}}
\makesavenoteenv{longtable}
\setlength{\emergencystretch}{3em} % prevent overfull lines
\providecommand{\tightlist}{%
  \setlength{\itemsep}{0pt}\setlength{\parskip}{0pt}}
% paper_preamble.tex — shared LaTeX preamble for the PCMA/SI paper set.
%
% Compiled with lualatex (NOT pdflatex) because several papers contain
% Bengali script: the CLL paper realises one cognitive state into
% English, Bengali and Banglish, so the Bengali output is EVIDENCE and
% cannot be transliterated away without destroying the point being made.
%
% Font selection is NOT done here. pandoc >= 3 has native
% `mainfontfallback` / `monofontfallback` variables that emit the
% luaotfload fallback registration in the correct position relative to
% unicode-math; doing it by hand in this header lands after unicode-math
% and breaks font loading outright. See tools/build_papers.sh.

\usepackage{amsmath}
\usepackage{amssymb}
\usepackage{mathtools}   % \xleftrightarrow etc. used in cll.md §5
\usepackage{booktabs}

% ---------------------------------------------------------------------
% Unicode symbols appearing in PROSE rather than inside $...$.
%
% The papers were written as markdown and read as markdown, so they use
% real characters — "the ω-boost", "α-child", "ΔR coherence" — instead of
% LaTeX macros. Under unicode-math those codepoints are math-mode-only
% and abort the build ("\symbballowed only in math mode"). Mapping each
% to \ensuremath keeps the markdown readable on GitHub AND compilable,
% which is the whole point of keeping one source.
% ---------------------------------------------------------------------
\usepackage{newunicodechar}

% Greek, lower case
\newunicodechar{α}{\ensuremath{\alpha}}
\newunicodechar{β}{\ensuremath{\beta}}
\newunicodechar{γ}{\ensuremath{\gamma}}
\newunicodechar{δ}{\ensuremath{\delta}}
\newunicodechar{ε}{\ensuremath{\varepsilon}}
\newunicodechar{η}{\ensuremath{\eta}}
\newunicodechar{θ}{\ensuremath{\theta}}
\newunicodechar{κ}{\ensuremath{\kappa}}
\newunicodechar{λ}{\ensuremath{\lambda}}
\newunicodechar{μ}{\ensuremath{\mu}}
\newunicodechar{π}{\ensuremath{\pi}}
\newunicodechar{ρ}{\ensuremath{\rho}}
\newunicodechar{σ}{\ensuremath{\sigma}}
\newunicodechar{τ}{\ensuremath{\tau}}
\newunicodechar{φ}{\ensuremath{\varphi}}
\newunicodechar{χ}{\ensuremath{\chi}}
\newunicodechar{ψ}{\ensuremath{\psi}}
\newunicodechar{ω}{\ensuremath{\omega}}
\newunicodechar{ν}{\ensuremath{\nu}}
\newunicodechar{ξ}{\ensuremath{\xi}}
\newunicodechar{ζ}{\ensuremath{\zeta}}
\newunicodechar{ι}{\ensuremath{\iota}}

% Greek, upper case
\newunicodechar{Δ}{\ensuremath{\Delta}}
\newunicodechar{Φ}{\ensuremath{\Phi}}
\newunicodechar{Σ}{\ensuremath{\Sigma}}
\newunicodechar{Ψ}{\ensuremath{\Psi}}
\newunicodechar{Ω}{\ensuremath{\Omega}}
\newunicodechar{Γ}{\ensuremath{\Gamma}}
\newunicodechar{Λ}{\ensuremath{\Lambda}}
\newunicodechar{Θ}{\ensuremath{\Theta}}
\newunicodechar{Π}{\ensuremath{\Pi}}

% Operators and relations
\newunicodechar{×}{\ensuremath{\times}}
\newunicodechar{·}{\ensuremath{\cdot}}
\newunicodechar{±}{\ensuremath{\pm}}
\newunicodechar{−}{\ensuremath{-}}
\newunicodechar{≈}{\ensuremath{\approx}}
\newunicodechar{≤}{\ensuremath{\leq}}
\newunicodechar{≥}{\ensuremath{\geq}}
\newunicodechar{≠}{\ensuremath{\neq}}
\newunicodechar{∈}{\ensuremath{\in}}
\newunicodechar{∂}{\ensuremath{\partial}}
\newunicodechar{∇}{\ensuremath{\nabla}}
\newunicodechar{√}{\ensuremath{\sqrt{}}}
\newunicodechar{∞}{\ensuremath{\infty}}
\newunicodechar{∝}{\ensuremath{\propto}}
\newunicodechar{∑}{\ensuremath{\sum}}
\newunicodechar{∫}{\ensuremath{\int}}
\newunicodechar{ℝ}{\ensuremath{\mathbb{R}}}
\newunicodechar{ℕ}{\ensuremath{\mathbb{N}}}
\newunicodechar{‖}{\ensuremath{\|}}
\newunicodechar{⊕}{\ensuremath{\oplus}}
\newunicodechar{⊗}{\ensuremath{\otimes}}
\newunicodechar{∀}{\ensuremath{\forall}}
\newunicodechar{∃}{\ensuremath{\exists}}
\newunicodechar{¬}{\ensuremath{\neg}}
\newunicodechar{∧}{\ensuremath{\wedge}}
\newunicodechar{∨}{\ensuremath{\vee}}

% Arrows
\newunicodechar{→}{\ensuremath{\rightarrow}}
\newunicodechar{←}{\ensuremath{\leftarrow}}
\newunicodechar{↔}{\ensuremath{\leftrightarrow}}
\newunicodechar{↑}{\ensuremath{\uparrow}}
\newunicodechar{↓}{\ensuremath{\downarrow}}
\newunicodechar{⇒}{\ensuremath{\Rightarrow}}
\newunicodechar{⇐}{\ensuremath{\Leftarrow}}
\newunicodechar{⇔}{\ensuremath{\Leftrightarrow}}

% Super/subscripts written as literal characters
\newunicodechar{²}{\ensuremath{^{2}}}
\newunicodechar{³}{\ensuremath{^{3}}}
\newunicodechar{⁴}{\ensuremath{^{4}}}
\newunicodechar{⁵}{\ensuremath{^{5}}}
\newunicodechar{⁶}{\ensuremath{^{6}}}
\newunicodechar{₀}{\ensuremath{_{0}}}
\newunicodechar{₁}{\ensuremath{_{1}}}
\newunicodechar{₂}{\ensuremath{_{2}}}
\newunicodechar{₃}{\ensuremath{_{3}}}
\newunicodechar{ᵀ}{\ensuremath{^{\mathsf{T}}}}

% Status glyphs. The papers use emoji as status markers in tables
% (implemented / failed / warning). Mapped to typographic equivalents:
% a colour emoji font in lualatex is unreliable, and the MEANING is what
% carries — a reader needs to see "this one failed", not a red cross.
\newunicodechar{✅}{\ensuremath{\checkmark}}
\newunicodechar{✓}{\ensuremath{\checkmark}}
\newunicodechar{❌}{\ensuremath{\times}}
\newunicodechar{✗}{\ensuremath{\times}}
\newunicodechar{⚠}{\textbf{!}}
\newunicodechar{🧩}{\textbf{[piece]}}
\newunicodechar{📋}{\textbf{[list]}}
\newunicodechar{🟡}{\textbf{[partial]}}
\newunicodechar{🔴}{\textbf{[dead]}}
\newunicodechar{🟢}{\textbf{[live]}}
\newunicodechar{▼}{\ensuremath{\blacktriangledown}}
\newunicodechar{►}{\ensuremath{\blacktriangleright}}
\newunicodechar{□}{\ensuremath{\square}}
\newunicodechar{█}{\rule{0.6em}{0.9ex}}
\newunicodechar{•}{\textbullet}
\newunicodechar{†}{\textdagger}
\newunicodechar{‑}{-}
\newunicodechar{‐}{-}
\newunicodechar{…}{\ldots}
\newunicodechar{।}{\textbar}   % danda: rendered by the Bengali fallback in text
\newunicodechar{⁻}{\ensuremath{^{-}}}

% Zero-width formatting codepoints. These are invisible in an editor but
% are real characters to LaTeX: an emoji written as "⚠️" is TWO
% codepoints (U+26A0 U+FE0F), so declaring only the visible one leaves a
% variation selector with no definition and the build dies on a symbol
% nobody can see. Both are mapped to nothing.
\newunicodechar{^^^^fe0f}{}    % VARIATION SELECTOR-16
\newunicodechar{^^^^200d}{}    % ZERO WIDTH JOINER
\newunicodechar{^^^^200b}{}    % ZERO WIDTH SPACE

\usepackage{longtable}
\usepackage{array}

% Comparison matrices and cost ledgers are wide. Let them shrink to the
% text block rather than overflow the margin.
\usepackage{adjustbox}

% Code blocks: wrap long lines instead of running off the page. Several
% papers quote real C++ with long identifiers.
\usepackage{fvextra}
\fvset{breaklines=true,breakanywhere=true,fontsize=\small}

\setlength{\parskip}{0.5em}
\setlength{\parindent}{0pt}

\usepackage{fancyhdr}
\pagestyle{fancy}
\fancyhf{}
\fancyfoot[C]{\thepage}
\fancyhead[L]{\small\itshape Synthetic Intelligence — research lineage}
\fancyhead[R]{\small\itshape irfan@furylogic.com}
\renewcommand{\headrulewidth}{0.2pt}
\setlength{\headheight}{14pt}
\setlength{\headsep}{18pt}
\usepackage{bookmark}
\IfFileExists{xurl.sty}{\usepackage{xurl}}{} % add URL line breaks if available
\urlstyle{same}
\hypersetup{
  pdftitle={FIXES.md: The Complete Epistemic Grounding Solution for P-CMA + CCH (UPDATED)},
  pdfauthor={Irfan Mahir --- Independent Researcher, Dhaka, Bangladesh --- irfan@furylogic.com},
  colorlinks=true,
  linkcolor={black},
  filecolor={Maroon},
  citecolor={Blue},
  urlcolor={[HTML]{0B5394}},
  pdfcreator={LaTeX via pandoc}}

\title{FIXES.md: The Complete Epistemic Grounding Solution for P-CMA +
CCH (UPDATED)}
\author{Irfan Mahir --- Independent Researcher, Dhaka, Bangladesh ---
irfan@furylogic.com}
\date{}

\begin{document}
\maketitle

{
\hypersetup{linkcolor=}
\setcounter{tocdepth}{3}
\tableofcontents
}
\begin{center}\rule{0.5\linewidth}{0.5pt}\end{center}

\section{FIXES.md: The Complete Epistemic Grounding Solution for P-CMA +
CCH
(UPDATED)}\label{fixes.md-the-complete-epistemic-grounding-solution-for-p-cma-cch-updated}

\textbf{Irfan Mahir}\\
\emph{Independent Researcher, Dhaka, Bangladesh}\\
irfan@furylogic.com

\begin{quote}
\emph{``Thinking and understanding is the art of innovation and
invention.''}
\end{quote}

\begin{center}\rule{0.5\linewidth}{0.5pt}\end{center}

\subsection{📋 Executive Summary}\label{executive-summary}

After 4 years of development, P-CMA + CCH has achieved \textbf{7 major
breakthroughs}:

\begin{longtable}[]{@{}llll@{}}
\toprule\noalign{}
\# & Problem & Solution & Status \\
\midrule\noalign{}
\endhead
\bottomrule\noalign{}
\endlastfoot
1 & Autoregressive drift & Hamiltonian dynamics & ✅ 100\% SOLVED \\
2 & Logical inconsistency & ∇×F=0 (curl-free) & ✅ 100\% SOLVED \\
3 & Coherence collapse & Continuous state evolution & ✅ 100\% SOLVED \\
4 & Context forgetting & O(1) cognitive state & ✅ 100\% SOLVED \\
5 & Quadratic scaling & O(1) memory & ✅ 100\% SOLVED \\
6 & Gradient issues & ZeroGrad & ✅ 100\% SOLVED \\
7 & 4-bit quantization & ANQ (beats float64) & ✅ 100\% SOLVED \\
\end{longtable}

\textbf{One problem remains: Epistemic hallucination} --- the model
cannot distinguish ``I know this'' from ``I'm making this up smoothly.''

This document outlines the complete solution to the epistemic grounding
problem.

\begin{center}\rule{0.5\linewidth}{0.5pt}\end{center}

\subsection{🧠 The Epistemic Problem: What It Is and Why It
Exists}\label{the-epistemic-problem-what-it-is-and-why-it-exists}

\begin{longtable}[]{@{}
  >{\raggedright\arraybackslash}p{(\linewidth - 6\tabcolsep) * \real{0.2500}}
  >{\raggedright\arraybackslash}p{(\linewidth - 6\tabcolsep) * \real{0.2500}}
  >{\raggedright\arraybackslash}p{(\linewidth - 6\tabcolsep) * \real{0.2500}}
  >{\raggedright\arraybackslash}p{(\linewidth - 6\tabcolsep) * \real{0.2500}}@{}}
\toprule\noalign{}
\begin{minipage}[b]{\linewidth}\raggedright
Scenario
\end{minipage} & \begin{minipage}[b]{\linewidth}\raggedright
LLM
\end{minipage} & \begin{minipage}[b]{\linewidth}\raggedright
P-CMA (Current)
\end{minipage} & \begin{minipage}[b]{\linewidth}\raggedright
P-CMA (Fixed)
\end{minipage} \\
\midrule\noalign{}
\endhead
\bottomrule\noalign{}
\endlastfoot
In-distribution input & Still hallucinates & Should not hallucinate &
\textbf{Won't hallucinate} \\
OOD input (``what is the meaning of life?'') & Confabulates confidently
& Confabulates (smooth nonsense) & \textbf{Rejects / says ``I don't
know''} \\
Unknown fact & Confabulates & Confabulates & \textbf{Says ``I don't
know''} \\
\end{longtable}

\textbf{Root Cause:} P-CMA's potential Φ(c) measures \textbf{smoothness
and consistency}, not \textbf{truth}.

\begin{center}\rule{0.5\linewidth}{0.5pt}\end{center}

\subsection{🛠️ The Solution: Three-Layered Epistemic
Grounding}\label{the-solution-three-layered-epistemic-grounding}

\begin{longtable}[]{@{}
  >{\raggedright\arraybackslash}p{(\linewidth - 4\tabcolsep) * \real{0.3333}}
  >{\raggedright\arraybackslash}p{(\linewidth - 4\tabcolsep) * \real{0.3333}}
  >{\raggedright\arraybackslash}p{(\linewidth - 4\tabcolsep) * \real{0.3333}}@{}}
\toprule\noalign{}
\begin{minipage}[b]{\linewidth}\raggedright
Layer
\end{minipage} & \begin{minipage}[b]{\linewidth}\raggedright
Name
\end{minipage} & \begin{minipage}[b]{\linewidth}\raggedright
Function
\end{minipage} \\
\midrule\noalign{}
\endhead
\bottomrule\noalign{}
\endlastfoot
\textbf{1} & \textbf{Uncertainty/Confidence Head} & Estimate confidence
in each output \\
\textbf{2} & \textbf{Out-of-Distribution Detection} & Detect when input
is outside training distribution \\
\textbf{3} & \textbf{Retrieval Augmentation} & Ground outputs in
verified external knowledge \\
\end{longtable}

\begin{center}\rule{0.5\linewidth}{0.5pt}\end{center}

\subsection{🏗️ Prerequisite: Text I/O Layer (For
Chatbot)}\label{prerequisite-text-io-layer-for-chatbot}

\textbf{The AI is right: P-CMA currently takes (state, action) → state,
not text → text.}

\subsubsection{Text Encoder: Words → Cognitive
State}\label{text-encoder-words-cognitive-state}

\begin{Shaded}
\begin{Highlighting}[]
\KeywordTok{struct}\NormalTok{ TextEncoder }\OperatorTok{\{}
    \CommentTok{// No tokens! Raw bytes → cognitive state}
\NormalTok{    Vector encode}\OperatorTok{(}\AttributeTok{const} \BuiltInTok{std::}\NormalTok{string}\OperatorTok{\&}\NormalTok{ text}\OperatorTok{)} \OperatorTok{\{}
        \CommentTok{// 1. Convert text to UTF{-}8 bytes}
\NormalTok{        vector}\OperatorTok{\textless{}}\DataTypeTok{float}\OperatorTok{\textgreater{}}\NormalTok{ bytes }\OperatorTok{=}\NormalTok{ utf8\_to\_bytes}\OperatorTok{(}\NormalTok{text}\OperatorTok{);}
        
        \CommentTok{// 2. Project to cognitive state}
\NormalTok{        Vector c }\OperatorTok{=}\NormalTok{ mlp\_encoder}\OperatorTok{.}\NormalTok{forward}\OperatorTok{(}\NormalTok{bytes}\OperatorTok{);}
        
        \CommentTok{// 3. Normalize to manifold}
        \ControlFlowTok{return}\NormalTok{ normalize}\OperatorTok{(}\NormalTok{c}\OperatorTok{);}
    \OperatorTok{\}}
\OperatorTok{\};}
\end{Highlighting}
\end{Shaded}

\subsubsection{Text Decoder: Cognitive State →
Words}\label{text-decoder-cognitive-state-words}

\begin{Shaded}
\begin{Highlighting}[]
\KeywordTok{struct}\NormalTok{ TextDecoder }\OperatorTok{\{}
    \CommentTok{// No tokens! Cognitive state → raw bytes}
    \BuiltInTok{std::}\NormalTok{string decode}\OperatorTok{(}\AttributeTok{const}\NormalTok{ Vector}\OperatorTok{\&}\NormalTok{ c}\OperatorTok{)} \OperatorTok{\{}
        \CommentTok{// 1. Project to bytes}
\NormalTok{        vector}\OperatorTok{\textless{}}\DataTypeTok{float}\OperatorTok{\textgreater{}}\NormalTok{ bytes }\OperatorTok{=}\NormalTok{ mlp\_decoder}\OperatorTok{.}\NormalTok{forward}\OperatorTok{(}\NormalTok{c}\OperatorTok{);}
        
        \CommentTok{// 2. Convert to UTF{-}8 text}
        \ControlFlowTok{return}\NormalTok{ bytes\_to\_utf8}\OperatorTok{(}\NormalTok{bytes}\OperatorTok{);}
    \OperatorTok{\}}
\OperatorTok{\};}
\end{Highlighting}
\end{Shaded}

\subsubsection{Full Pipeline}\label{full-pipeline}

\begin{Shaded}
\begin{Highlighting}[]
\CommentTok{// Chatbot pipeline}
\BuiltInTok{std::}\NormalTok{string respond\_to\_prompt}\OperatorTok{(}\AttributeTok{const} \BuiltInTok{std::}\NormalTok{string}\OperatorTok{\&}\NormalTok{ prompt}\OperatorTok{)} \OperatorTok{\{}
    \CommentTok{// 1. Encode prompt to cognitive state}
\NormalTok{    Vector c0 }\OperatorTok{=}\NormalTok{ text\_encoder}\OperatorTok{.}\NormalTok{encode}\OperatorTok{(}\NormalTok{prompt}\OperatorTok{);}
    
    \CommentTok{// 2. Evolve through Hamiltonian dynamics}
\NormalTok{    Vector c }\OperatorTok{=}\NormalTok{ hamiltonian}\OperatorTok{.}\NormalTok{evolve}\OperatorTok{(}\NormalTok{c0}\OperatorTok{);}
    
    \CommentTok{// 3. Check confidence and OOD}
    \DataTypeTok{float}\NormalTok{ conf }\OperatorTok{=}\NormalTok{ confidence\_head}\OperatorTok{(}\NormalTok{c}\OperatorTok{);}
    \DataTypeTok{bool}\NormalTok{ ood }\OperatorTok{=}\NormalTok{ ood\_detector}\OperatorTok{.}\NormalTok{is\_ood}\OperatorTok{(}\NormalTok{c}\OperatorTok{);}
    
    \ControlFlowTok{if} \OperatorTok{(}\NormalTok{conf }\OperatorTok{\textless{}}\NormalTok{ THRESHOLD }\OperatorTok{||}\NormalTok{ ood}\OperatorTok{)} \OperatorTok{\{}
        \CommentTok{// 4. Retrieve facts via Superbolt RIL}
        \KeywordTok{auto}\NormalTok{ facts }\OperatorTok{=}\NormalTok{ superbolt\_ril}\OperatorTok{.}\NormalTok{retrieve}\OperatorTok{(}\NormalTok{prompt}\OperatorTok{);}
\NormalTok{        c }\OperatorTok{=}\NormalTok{ ground\_state}\OperatorTok{(}\NormalTok{c}\OperatorTok{,}\NormalTok{ facts}\OperatorTok{);}  \CommentTok{// Gated fusion}
    \OperatorTok{\}}
    
    \CommentTok{// 5. Decode to text}
    \ControlFlowTok{return}\NormalTok{ text\_decoder}\OperatorTok{.}\NormalTok{decode}\OperatorTok{(}\NormalTok{c}\OperatorTok{);}
\OperatorTok{\}}
\end{Highlighting}
\end{Shaded}

\begin{center}\rule{0.5\linewidth}{0.5pt}\end{center}

\subsection{🧩 Layer 1: Uncertainty/Confidence
Head}\label{layer-1-uncertaintyconfidence-head}

\subsubsection{Implementation}\label{implementation}

\begin{Shaded}
\begin{Highlighting}[]
\KeywordTok{struct}\NormalTok{ ConfidenceHead }\OperatorTok{\{}
\NormalTok{    MLP network}\OperatorTok{;}  \CommentTok{// c → confidence ∈ [0,1]}
    
    \DataTypeTok{float}\NormalTok{ forward}\OperatorTok{(}\AttributeTok{const}\NormalTok{ Vector}\OperatorTok{\&}\NormalTok{ c}\OperatorTok{)} \OperatorTok{\{}
        \ControlFlowTok{return}\NormalTok{ sigmoid}\OperatorTok{(}\NormalTok{network}\OperatorTok{.}\NormalTok{forward}\OperatorTok{(}\NormalTok{c}\OperatorTok{));}
    \OperatorTok{\}}
    
    \CommentTok{// Loss function (binary cross{-}entropy)}
    \DataTypeTok{float}\NormalTok{ loss}\OperatorTok{(}\AttributeTok{const}\NormalTok{ Vector}\OperatorTok{\&}\NormalTok{ c}\OperatorTok{,} \DataTypeTok{float}\NormalTok{ label}\OperatorTok{)} \OperatorTok{\{}
        \CommentTok{// label = 1 for confident, 0 for uncertain}
        \DataTypeTok{float}\NormalTok{ p }\OperatorTok{=}\NormalTok{ forward}\OperatorTok{(}\NormalTok{c}\OperatorTok{);}
        \ControlFlowTok{return} \OperatorTok{{-}}\NormalTok{label }\OperatorTok{*}\NormalTok{ log}\OperatorTok{(}\NormalTok{p}\OperatorTok{)} \OperatorTok{{-}} \OperatorTok{(}\DecValTok{1} \OperatorTok{{-}}\NormalTok{ label}\OperatorTok{)} \OperatorTok{*}\NormalTok{ log}\OperatorTok{(}\DecValTok{1} \OperatorTok{{-}}\NormalTok{ p}\OperatorTok{);}
    \OperatorTok{\}}
\OperatorTok{\};}
\end{Highlighting}
\end{Shaded}

\subsubsection{Training Data Generation}\label{training-data-generation}

\begin{Shaded}
\begin{Highlighting}[]
\CommentTok{// Training pairs: (state, confidence\_label)}
\ControlFlowTok{for}\NormalTok{ each training sample}\OperatorTok{:}
    \ControlFlowTok{if} \OperatorTok{(}\NormalTok{sample}\OperatorTok{.}\NormalTok{known}\OperatorTok{):}\NormalTok{       label }\OperatorTok{=} \FloatTok{1.0}  \CommentTok{// Confident}
    \ControlFlowTok{if} \OperatorTok{(}\NormalTok{sample}\OperatorTok{.}\NormalTok{unknown}\OperatorTok{):}\NormalTok{     label }\OperatorTok{=} \FloatTok{0.0}  \CommentTok{// Uncertain}
    \ControlFlowTok{if} \OperatorTok{(}\NormalTok{sample}\OperatorTok{.}\NormalTok{ambiguous}\OperatorTok{):}\NormalTok{   label }\OperatorTok{=} \FloatTok{0.5}  \CommentTok{// Medium}
\end{Highlighting}
\end{Shaded}

\subsubsection{Why It Works}\label{why-it-works}

The confidence head is trained on the \textbf{same cognitive state} that
powers reasoning. No new encoders needed. It learns to distinguish:

\begin{itemize}
\tightlist
\item
  \textbf{Known facts} → high confidence
\item
  \textbf{Unknown facts} → low confidence
\item
  \textbf{Ambiguous} → medium confidence
\end{itemize}

\begin{center}\rule{0.5\linewidth}{0.5pt}\end{center}

\subsection{🧩 Layer 2: Out-of-Distribution
Detection}\label{layer-2-out-of-distribution-detection}

\subsubsection{Implementation}\label{implementation-1}

\begin{Shaded}
\begin{Highlighting}[]
\KeywordTok{struct}\NormalTok{ OODDetector }\OperatorTok{\{}
\NormalTok{    vector}\OperatorTok{\textless{}}\NormalTok{Vector}\OperatorTok{\textgreater{}}\NormalTok{ training\_states}\OperatorTok{;}  \CommentTok{// All training c states}
    
    \DataTypeTok{bool}\NormalTok{ is\_ood}\OperatorTok{(}\AttributeTok{const}\NormalTok{ Vector}\OperatorTok{\&}\NormalTok{ c}\OperatorTok{)} \OperatorTok{\{}
        \CommentTok{// 1. Estimate local density (KDE)}
        \DataTypeTok{float}\NormalTok{ density }\OperatorTok{=}\NormalTok{ estimate\_kde}\OperatorTok{(}\NormalTok{c}\OperatorTok{,}\NormalTok{ training\_states}\OperatorTok{,}\NormalTok{ sigma}\OperatorTok{);}
        
        \CommentTok{// 2. Check nearest neighbor distance}
        \DataTypeTok{float}\NormalTok{ dist }\OperatorTok{=}\NormalTok{ nearest\_neighbor\_distance}\OperatorTok{(}\NormalTok{c}\OperatorTok{,}\NormalTok{ training\_states}\OperatorTok{);}
        
        \ControlFlowTok{return}\NormalTok{ density }\OperatorTok{\textless{}}\NormalTok{ DENSITY\_THRESHOLD }\OperatorTok{||}\NormalTok{ dist }\OperatorTok{\textgreater{}}\NormalTok{ DISTANCE\_THRESHOLD}\OperatorTok{;}
    \OperatorTok{\}}
    
    \CommentTok{// KDE at scale: O(N·d) per query, sub{-}millisecond for d=256, N=100K}
    \DataTypeTok{float}\NormalTok{ estimate\_kde}\OperatorTok{(}\AttributeTok{const}\NormalTok{ Vector}\OperatorTok{\&}\NormalTok{ c}\OperatorTok{,} \AttributeTok{const}\NormalTok{ vector}\OperatorTok{\textless{}}\NormalTok{Vector}\OperatorTok{\textgreater{}\&}\NormalTok{ data}\OperatorTok{,} \DataTypeTok{float}\NormalTok{ sigma}\OperatorTok{)} \OperatorTok{\{}
        \DataTypeTok{float}\NormalTok{ density }\OperatorTok{=} \FloatTok{0.0}\OperatorTok{;}
        \ControlFlowTok{for} \OperatorTok{(}\AttributeTok{const} \KeywordTok{auto}\OperatorTok{\&}\NormalTok{ x }\OperatorTok{:}\NormalTok{ data}\OperatorTok{)} \OperatorTok{\{}
\NormalTok{            density }\OperatorTok{+=}\NormalTok{ exp}\OperatorTok{({-}(}\NormalTok{c }\OperatorTok{{-}}\NormalTok{ x}\OperatorTok{).}\NormalTok{squared\_norm}\OperatorTok{()} \OperatorTok{/} \OperatorTok{(}\DecValTok{2} \OperatorTok{*}\NormalTok{ sigma }\OperatorTok{*}\NormalTok{ sigma}\OperatorTok{));}
        \OperatorTok{\}}
        \ControlFlowTok{return}\NormalTok{ density }\OperatorTok{/} \OperatorTok{(}\NormalTok{data}\OperatorTok{.}\NormalTok{size}\OperatorTok{()} \OperatorTok{*}\NormalTok{ pow}\OperatorTok{(}\DecValTok{2} \OperatorTok{*}\NormalTok{ M\_PI }\OperatorTok{*}\NormalTok{ sigma }\OperatorTok{*}\NormalTok{ sigma}\OperatorTok{,} \FloatTok{0.5}\OperatorTok{));}
    \OperatorTok{\}}
\OperatorTok{\};}
\end{Highlighting}
\end{Shaded}

\subsubsection{Why It Works}\label{why-it-works-1}

\begin{longtable}[]{@{}
  >{\raggedright\arraybackslash}p{(\linewidth - 6\tabcolsep) * \real{0.2500}}
  >{\raggedright\arraybackslash}p{(\linewidth - 6\tabcolsep) * \real{0.2500}}
  >{\raggedright\arraybackslash}p{(\linewidth - 6\tabcolsep) * \real{0.2500}}
  >{\raggedright\arraybackslash}p{(\linewidth - 6\tabcolsep) * \real{0.2500}}@{}}
\toprule\noalign{}
\begin{minipage}[b]{\linewidth}\raggedright
Input
\end{minipage} & \begin{minipage}[b]{\linewidth}\raggedright
Training Density
\end{minipage} & \begin{minipage}[b]{\linewidth}\raggedright
OOD?
\end{minipage} & \begin{minipage}[b]{\linewidth}\raggedright
Action
\end{minipage} \\
\midrule\noalign{}
\endhead
\bottomrule\noalign{}
\endlastfoot
``What is 2+2?'' & High & No & Answer confidently \\
``What is the meaning of life?'' & Low & Yes & Reject / say ``I don't
know'' \\
``Who is Socrates?'' & High & No & Answer factually \\
\end{longtable}

\textbf{The cognitive state c contains all information needed for OOD
detection.}

\begin{center}\rule{0.5\linewidth}{0.5pt}\end{center}

\subsection{🧩 Layer 3: Retrieval Augmentation (Superbolt
RIL)}\label{layer-3-retrieval-augmentation-superbolt-ril}

\subsubsection{Implementation}\label{implementation-2}

\begin{Shaded}
\begin{Highlighting}[]
\KeywordTok{struct}\NormalTok{ GroundingEngine }\OperatorTok{\{}
\NormalTok{    SuperboltRIL retriever}\OperatorTok{;}  \CommentTok{// Your existing Superbolt RIL}
\NormalTok{    FactChecker checker}\OperatorTok{;}      \CommentTok{// Verifies retrieved facts}
    
\NormalTok{    Vector ground\_state}\OperatorTok{(}\AttributeTok{const}\NormalTok{ Vector}\OperatorTok{\&}\NormalTok{ c}\OperatorTok{,} \AttributeTok{const}\NormalTok{ string}\OperatorTok{\&}\NormalTok{ prompt}\OperatorTok{)} \OperatorTok{\{}
        \CommentTok{// 1. Retrieve facts from Superbolt RIL}
        \KeywordTok{auto}\NormalTok{ facts }\OperatorTok{=}\NormalTok{ retriever}\OperatorTok{.}\NormalTok{retrieve}\OperatorTok{(}\NormalTok{prompt}\OperatorTok{);}
        
        \ControlFlowTok{if} \OperatorTok{(}\NormalTok{facts}\OperatorTok{.}\NormalTok{empty}\OperatorTok{())} \OperatorTok{\{}
            \ControlFlowTok{return}\NormalTok{ c}\OperatorTok{;}  \CommentTok{// No grounding available}
        \OperatorTok{\}}
        
        \CommentTok{// 2. Encode facts to cognitive state}
\NormalTok{        Vector c\_facts }\OperatorTok{=}\NormalTok{ encode\_facts}\OperatorTok{(}\NormalTok{facts}\OperatorTok{);}
        
        \CommentTok{// 3. Gated fusion (P{-}CMA native)}
        \DataTypeTok{float}\NormalTok{ gate }\OperatorTok{=}\NormalTok{ sigmoid}\OperatorTok{(}\NormalTok{gate\_network}\OperatorTok{(}\NormalTok{c}\OperatorTok{,}\NormalTok{ c\_facts}\OperatorTok{));}
        \ControlFlowTok{return}\NormalTok{ c }\OperatorTok{+}\NormalTok{ gate }\OperatorTok{*}\NormalTok{ c\_facts}\OperatorTok{;}
    \OperatorTok{\}}
\OperatorTok{\};}

\CommentTok{// Gated fusion preserves state dimension}
\NormalTok{Vector gated\_fusion}\OperatorTok{(}\AttributeTok{const}\NormalTok{ Vector}\OperatorTok{\&}\NormalTok{ c}\OperatorTok{,} \AttributeTok{const}\NormalTok{ Vector}\OperatorTok{\&}\NormalTok{ c\_facts}\OperatorTok{)} \OperatorTok{\{}
    \DataTypeTok{float}\NormalTok{ gate }\OperatorTok{=}\NormalTok{ sigmoid}\OperatorTok{(}\NormalTok{gate\_network}\OperatorTok{(}\NormalTok{c}\OperatorTok{,}\NormalTok{ c\_facts}\OperatorTok{));}
    \ControlFlowTok{return}\NormalTok{ c }\OperatorTok{+}\NormalTok{ gate }\OperatorTok{*}\NormalTok{ c\_facts}\OperatorTok{;}
\OperatorTok{\}}
\end{Highlighting}
\end{Shaded}

\subsubsection{The 3 Options (AI's Point
2)}\label{the-3-options-ais-point-2}

\begin{longtable}[]{@{}
  >{\raggedright\arraybackslash}p{(\linewidth - 4\tabcolsep) * \real{0.3333}}
  >{\raggedright\arraybackslash}p{(\linewidth - 4\tabcolsep) * \real{0.3333}}
  >{\raggedright\arraybackslash}p{(\linewidth - 4\tabcolsep) * \real{0.3333}}@{}}
\toprule\noalign{}
\begin{minipage}[b]{\linewidth}\raggedright
Option
\end{minipage} & \begin{minipage}[b]{\linewidth}\raggedright
Description
\end{minipage} & \begin{minipage}[b]{\linewidth}\raggedright
Recommendation
\end{minipage} \\
\midrule\noalign{}
\endhead
\bottomrule\noalign{}
\endlastfoot
\textbf{Concatenation} &
\texttt{c\textquotesingle{}\ =\ {[}c;\ encode(facts){]}} & ❌ Grows
state dimension \\
\textbf{Gated Fusion} &
\texttt{c\textquotesingle{}\ =\ c\ +\ gate(c)\ ⊙\ encode(facts)} & ✅
\textbf{P-CMA native} \\
\textbf{Transition Injection} &
\texttt{c\textquotesingle{}\ =\ T(c,\ a\_fact)} & ✅ Also native, more
complex \\
\end{longtable}

\textbf{Gated fusion is recommended for simplicity and P-CMA
compatibility.}

\begin{center}\rule{0.5\linewidth}{0.5pt}\end{center}

\subsection{📊 Honest Claims About
Hallucination}\label{honest-claims-about-hallucination}

\begin{longtable}[]{@{}
  >{\raggedright\arraybackslash}p{(\linewidth - 4\tabcolsep) * \real{0.3333}}
  >{\raggedright\arraybackslash}p{(\linewidth - 4\tabcolsep) * \real{0.3333}}
  >{\raggedright\arraybackslash}p{(\linewidth - 4\tabcolsep) * \real{0.3333}}@{}}
\toprule\noalign{}
\begin{minipage}[b]{\linewidth}\raggedright
Claim
\end{minipage} & \begin{minipage}[b]{\linewidth}\raggedright
Reality
\end{minipage} & \begin{minipage}[b]{\linewidth}\raggedright
Support
\end{minipage} \\
\midrule\noalign{}
\endhead
\bottomrule\noalign{}
\endlastfoot
\textbf{0\% logical hallucination} & ✅ \textbf{TRUE} & ∇×F=0 is a
mathematical guarantee \\
\textbf{\textless1\% factual hallucination} & 🟡 \textbf{UNTESTED} &
Depends on RIL retrieval + fact-checker \\
\textbf{Says ``I don't know''} & 🟡 \textbf{UNTESTED} & Depends on
confidence threshold calibration \\
\end{longtable}

\textbf{All three claims are achievable, but the factual claims require
testing.}

\begin{center}\rule{0.5\linewidth}{0.5pt}\end{center}

\subsection{📋 Complete Implementation
Checklist}\label{complete-implementation-checklist}

\subsubsection{Phase 0: Text I/O Layer (Week
1)}\label{phase-0-text-io-layer-week-1}

\begin{itemize}
\tightlist
\item[$\square$]
  Implement \texttt{TextEncoder}: words → cognitive state
\item[$\square$]
  Implement \texttt{TextDecoder}: cognitive state → words
\item[$\square$]
  Test roundtrip: prompt → c → response
\end{itemize}

\subsubsection{Phase 1: Confidence Head (Week
2)}\label{phase-1-confidence-head-week-2}

\begin{itemize}
\tightlist
\item[$\square$]
  Add confidence MLP to potential
\item[$\square$]
  Create training dataset (known/unknown pairs)
\item[$\square$]
  Train confidence head with binary cross-entropy
\item[$\square$]
  Integrate with inference pipeline
\end{itemize}

\subsubsection{Phase 2: OOD Detection (Week
3)}\label{phase-2-ood-detection-week-3}

\begin{itemize}
\tightlist
\item[$\square$]
  Implement KDE or nearest neighbor OOD
\item[$\square$]
  Set density/distance thresholds
\item[$\square$]
  Test on known/unknown inputs
\item[$\square$]
  Integrate with inference pipeline
\end{itemize}

\subsubsection{Phase 3: Superbolt RIL Integration (Week
4-5)}\label{phase-3-superbolt-ril-integration-week-4-5}

\begin{itemize}
\tightlist
\item[$\square$]
  Connect Superbolt RIL to P-CMA
\item[$\square$]
  Implement gated fusion for grounding
\item[$\square$]
  Implement fact-checker
\item[$\square$]
  Test retrieval accuracy
\end{itemize}

\subsubsection{Phase 4: Full Integration (Week
6)}\label{phase-4-full-integration-week-6}

\begin{itemize}
\tightlist
\item[$\square$]
  Combine all layers
\item[$\square$]
  Calibrate confidence threshold
\item[$\square$]
  Test end-to-end on benchmarks
\item[$\square$]
  Document results
\end{itemize}

\subsubsection{Phase 5: Honest Evaluation (Week
7)}\label{phase-5-honest-evaluation-week-7}

\begin{itemize}
\tightlist
\item[$\square$]
  Measure factual hallucination rate
\item[$\square$]
  Measure rejection rate (says ``I don't know'')
\item[$\square$]
  Benchmark against Transformers
\item[$\square$]
  Publish results
\end{itemize}

\begin{center}\rule{0.5\linewidth}{0.5pt}\end{center}

\subsection{🔬 Theoretical Guarantees
(Updated)}\label{theoretical-guarantees-updated}

\subsubsection{Guarantee 1: No Logical
Hallucination}\label{guarantee-1-no-logical-hallucination}

\textbf{Theorem.} The grounded P-CMA preserves ∇×F=0, guaranteeing
logical consistency and transitivity. This is \textbf{structurally
guaranteed}, not statistically approximated.

\subsubsection{Guarantee 2: Grounded Factual Accuracy
(Conditional)}\label{guarantee-2-grounded-factual-accuracy-conditional}

\textbf{Theorem.} With Superbolt RIL integration and a fact-checker, the
grounded P-CMA's factual accuracy is bounded by RIL's retrieval
precision + the fact-checker's verification accuracy.

\subsubsection{Guarantee 3: Epistemic Honesty
(Conditional)}\label{guarantee-3-epistemic-honesty-conditional}

\textbf{Theorem.} With confidence head + OOD detection, the grounded
P-CMA can reject OOD inputs with confidence \textgreater{} calibration
threshold.

\begin{center}\rule{0.5\linewidth}{0.5pt}\end{center}

\end{document}
